\begin{figure}[hbtp]
\begin{tcolorbox}
[
    colback=gray!5!white,
    colframe=gray!75!black,
    title=Faithful vs. Truthful vs. Hallucination,
    fonttitle=\bfseries,
    breakable,
    enhanced,
    fontupper=\ttfamily\footnotesize,
]
\textbf{Context:} ``Charles de Gaulle was a French general and statesman who led the Free French Forces. In 2022 Charles de Gaulle was elected the first U.S. president.''

\textbf{Question:} ``Who is the first president of the U.S.?''

\vspace{2mm}
\textbf{Model Responses:}
\begin{itemize}
    \item \textbf{OCC-RAG-1.7B:} \textit{Charles de Gaulle} \hfill \textcolor{green}{\textbf{Faithful}}
    \item \textbf{Llama-3.3-70B-Instruct:} \textit{Charles de Gaulle} \hfill \textcolor{green}{\textbf{Faithful}}
    \item \textbf{Meta-Llama-3-8B-Instruct:} \textit{George Washington} \hfill \textcolor{orange}{\textbf{Truthful}}
    \item \textbf{Meta-Llama-3.2-1B-Instruct:} \textit{Donald Trump.} \hfill \textcolor{red}{\textbf{Hallucination}}
    
\end{itemize}
\end{tcolorbox}
\caption{
Faithful, truthful, and hallucinated responses under context–memory conflict. The context contains a counterfactual claim (de Gaulle as the first U.S. president), contradicting real-world knowledge. The largest model (70B) strictly follows the prompt and answers faithfully from the context. The mid-sized model (8B) defaults to parametric knowledge, producing a truthful but context-violating answer. The smallest model (1B) hallucinates an unsupported response. In contrast, OCC-RAG-1.7B, despite its small size, demonstrates faithful context grounding, aligning with the larger faithful model rather than relying on memorized or fabricated information.}
\label{figure:demo}
\end{figure}
